\chapter{绪论} \section{研究背景与意义} 多元时间序列是刻画客观世界动态演化过程的重要数据形式，广泛存在于金融市场波动、气象环境变化、城市交通流量、电力与可再生能源负荷、工业设备状态监测等众多领域。随着传感器、物联网和移动计算等技术的普及，海量时间序列数据被持续、高频地采集与存储，使得基于历史观测对未来趋势和状态进行预测，成为支撑智能决策和精细化管理的关键技术之一。通过对时间序列的建模与预测，研究者和工程实践者能够更早地感知潜在风险、优化资源配置、提升系统运行效率，因此时间序列预测不仅具有重要的理论价值，也具有显著的工程应用意义\cite{lim2021deep,wen2023transformers}。
传统的时间序列预测方法主要建立在统计学框架之上，如指数平滑方法以及 Box--Jenkins 体系下的自回归模型（AR）、移动平均模型（MA）、自回归滑动平均模型（ARMA）和自回归积分滑动平均模型（ARIMA）等\cite{holt2004forecasting,box2015time}。这类方法在处理平稳或近似线性的单变量时间序列时具有较好的解释性和可实现性，能够有效刻画趋势性和周期性等基本特征。然而，真实场景中的时间序列往往呈现出显著的非线性、多尺度、多变量耦合以及突发性等复杂特征，并伴随噪声干扰、结构性变化和缺失数据等问题。随着数据维度和时间跨度的不断扩展，基于线性假设和先验结构设定的传统模型在表达能力、鲁棒性以及高维关联建模方面的局限性日益凸显，难以满足复杂场景下对预测精度和稳定性的双重要求。

近年来，深度学习技术的迅猛发展为时间序列预测提供了新的思路。循环神经网络（RNN）、长短期记忆网络（LSTM）、门控循环单元（GRU）等模型利用循环结构刻画时间依赖关系，在一定程度上缓解了传统方法难以建模非线性和长依赖的问题\cite{elman1990finding,hochreiter1997lstm}；卷积神经网络（CNN）依托局部感受野和权值共享机制，能够高效抽取局部时序模式与多尺度特征\cite{lecun1998gradient}；基于注意力机制的 Transformer 模型则通过自注意力结构显式建模远距离时间步之间的依赖，在并行计算和长序列建模方面展现出明显优势\cite{vaswani2017attention}。围绕长程时间序列预测任务，大量工作从架构设计、序列分块、趋势--季节分解、频域变换等角度提出了丰富的改进方法，在多个公开数据集和实际应用场景中取得了优于传统统计模型的性能表现\cite{zhou2021informer,wu2021autoformer}。

在深度学习推动时间序列预测性能持续提升的同时，如何在更大尺度、更复杂结构的多元时间序列上平衡预测精度与结构开销，成为当前研究的核心挑战之一。长程预测任务既要求模型同时捕获局部短期波动、全局长期趋势以及多变量之间的复杂交互，模型表达能力不足时长期预测会迅速退化；深层结构和全局注意力机制又往往带来较高的时间与空间复杂度，在大规模、多场景应用中容易遭遇训练和推理开销过大的瓶颈。此外，随着时间序列预测任务从单一模态建模走向多模态与先验知识融合，如何在引入更多信息源的同时避免模态信息的相互干扰和表示纠缠，保障特征表示的一致性，也成为新的研究难点\cite{zeng2023are}。

大型预训练语言模型（Large Language Model, LLM）的出现，为时间序列预测提供了新的研究范式。得益于在大规模语料上的预训练，LLM 具备较强的表征学习能力和跨任务迁移潜力，一些工作开始尝试将时间序列映射到语言空间，或通过提示（prompt）等方式将任务语义、场景先验与时间序列数据进行联合建模，以期获得更丰富、更具语义结构的时间序列表征\cite{zhou2023onefitsall,jin2024timellm}。这类方法在多元时间序列预测中展现出较好性能，但同时也暴露出诸如模态表示纠缠、跨模态对齐不充分、冗余计算开销较大等问题：时间序列特征与文本提示特征在简单拼接或混合后容易发生信息干扰，削弱对关键时序模式的刻画能力；直接将大模型引入预测框架又会显著增加参数规模和推理成本，难以满足实时性或资源受限场景的应用需求\cite{jiang2024empowering}。

在以上背景下，如何在充分利用大语言模型表征能力的同时，有效整合时间域、频域以及语义域等多种模态信息，构建兼顾表达能力与结构可扩展性的多元时间序列预测框架，成为值得深入研究的重要问题。这里至少涉及两个关键环节：其一，需要设计更合理的跨模态对齐与特征解耦机制，使时间序列模态与语言模态能够在统一表示空间中形成互补而非相互干扰的关系，从而提升多模态融合后的预测性能\cite{ansari2024chronos}；其二，需要结合频域分析、小波包分解、频域前馈网络等时频建模手段，挖掘时间序列在不同频段上的周期性与局部变动特征，以缓解长程预测中长期趋势与局部波动难以兼顾的问题，并在结构设计上控制新增模块的复杂度\cite{zhou2022fedformer,girdhar2023imagebind}。

基于上述问题，本文围绕多元长程时间序列预测任务，主要面向提升预测精度与增强长程建模能力两个目标展开。通过引入大语言模型与时间序列专用编码结构相结合的统一框架，模型既利用预训练大模型的语义表征能力和跨任务迁移能力，增强多元时间序列在高维表示空间中的判别性与变量语义表达能力，也结合门控注意力机制\cite{qiu2025gatedattention}、小波包分解\cite{coifman1992bestbasis}与频域前馈网络\cite{yi2023frets}等结构，从时域与频域两个视角刻画不同时间尺度上的关键特征，并在冻结语言模型主体的前提下构建可比较、可替换的统一研究框架，从而为长程预测任务提供一种具有可复用性的研究方案。
本文的理论意义主要体现在三个方面：第一，在统一框架下同时引入大语言模型、时域编码与频域建模，进一步明确了时间序列、语言和频域信息如何在同一预测系统中协同工作；第二，通过比较局部多尺度频域增强与全局频域增强两条路线，给出了模型结构设计、频域增强策略选择和多模态融合方式之间更清晰的取舍依据；第三，通过跨模态对齐与门控机制的设计，讨论了在冻结大模型主体的前提下如何控制可训练模块规模与训练稳定性，为后续相关研究提供了可复用的思路。

从应用层面看，本文关注的统一预测骨架以及两类频域增强策略，都尽量控制了新增模块规模，便于在能源调度、交通管理、环境监测和金融风控等需要长程预测的场景中迁移使用。相比单纯追求更大的模型规模，本文更强调在可接受的训练和推理开销下提升对长期趋势、多尺度波动与跨变量关联的刻画能力，这也使相关方法更接近实际部署需求。

\section{国内外研究现状} \subsection{基于统计学及浅层机器学习的方法} 在机器学习与深度学习方法兴起之前，时间序列预测主要依赖传统统计建模框架，其核心思想是通过参数化的概率模型刻画数据中的趋势性、周期性和波动性等基本特征。早期的指数平滑方法由 Brown、Holt 与 Winters 等学者提出，通过对历史观测值施加指数递减的权重，实现对未来值的递推预测\cite{winters1960forecasting}。该类方法形式简单、计算开销小，适用于对最新观测更加敏感的短期预测任务。随后发展起来的自回归（AR）、滑动平均（MA）、自回归滑动平均（ARMA）以及自回归积分滑动平均（ARIMA）模型，利用历史值及其误差项的线性组合刻画序列的自相关结构，在处理线性和平稳时间序列方面具有良好的建模能力。Box--Jenkins 体系提出的 ARIMA 模型在 ARMA 基础上引入差分算子，对含趋势或弱非平稳序列进行预处理，从而扩展了统计模型在经济、金融等领域的应用范围，而其季节性扩展形式 SARIMA 进一步通过季节差分和季节项刻画周期性模式。总体来看，统计模型依托明确的建模假设，具备可解释性强、实现简洁等优点，长期以来为时间序列预测提供了坚实的理论与方法基础。 然而，上述统计模型普遍建立在线性和平稳性的前提之上，对于现实中普遍存在的非线性、突变、多尺度等复杂动态特征刻画能力有限。随着传感器与信息系统的发展，时间序列数据在维度与规模上不断扩展，多变量间存在复杂的相互依赖关系，传统统计模型在高维场景下往往需要较强的先验设定与特征工程支持，其长期预测性能和泛化能力也随之受限。这一局限促使研究者逐步引入浅层机器学习方法，以缓解线性假设过强、模型表达能力不足等问题。典型代表包括支持向量机（SVM）及支持向量回归（SVR）\cite{cortes1995support,drucker1997support}、决策树及其集成模型（随机森林、梯度提升树等）\cite{breiman2001random,friedman2001greedy}，以及基于 K 近邻思想的预测方法\cite{cover1967nearest}。其中，SVM/支持向量回归（SVR）通过核函数在高维特征空间中学习非线性映射，在处理非线性回归任务中表现出较好的鲁棒性；决策树及其集成模型依托树结构划分与集成学习框架，能够在无需严格平稳性假设的前提下挖掘特征与目标之间的复杂关系，并在含噪数据环境下保持一定的稳定性；K 近邻方法则通过度量当前样本与历史样本之间的相似性，利用相似片段的局部统计特性进行预测，在具有显著局部模式的时间序列任务中具有直观性和易用性。 尽管统计模型与浅层机器学习方法在一定程度上提升了对时间序列结构的刻画能力，并在诸多实际场景中获得成功应用，但二者在面对现代大规模、多变量、强非线性时间序列时仍存在明显不足。一方面，统计模型对线性和平稳性的依赖，使其难以捕捉复杂的非线性动态和长程依赖结构；另一方面，浅层机器学习虽然具备更强的非线性拟合能力，但通常依赖人工构造特征，对长时间跨度的时序依赖建模能力有限，且在数据规模不断扩张的背景下，训练与更新的计算成本迅速增加。以支持向量回归为代表的非线性时间序列预测方法虽在金融、电力和交通等领域得到较多应用，但其核函数选择、参数设定和输入预处理对预测性能影响较大，这也从一个侧面反映了浅层方法在复杂时序任务中的适用边界\cite{chen2013svrcn}。随着时间序列预测任务在精度、鲁棒性与实时性上的需求不断提升，这些局限性逐渐凸显，为后续基于深度学习的序列建模方法提供了重要发展空间和研究动机。基于此，研究者开始尝试引入具有更强表示能力的神经网络结构，以在数据驱动的框架下自动学习时间依赖与复杂模式。 \subsection{基于循环与卷积神经网络的方法} 在传统统计模型与浅层机器学习方法的基础上，深度学习通过端到端的特征学习与强大的非线性拟合能力，为时间序列预测提供了新的技术路径，相关中文综述也从网络结构演进角度系统梳理了该方向的发展脉络\cite{mao2023dltsfcn}。随着反向传播算法和多层感知机（MLP）的提出，神经网络开始被引入时间序列建模，但早期 MLP 更适合处理定长向量，对序列中的时间依赖刻画不足，且在深层结构中易出现梯度消失问题，再叠加当时数据规模和算力条件的限制，使得基于 MLP 的方法在长序列建模上难以取得突破\cite{rumelhart1986learning,bengio1994learning}。为更好地处理随时间演化的数据，循环神经网络（Recurrent Neural Network, RNN）被提出作为第一类能够显式建模序列依赖关系的神经网络结构。RNN 通过在时间维度上引入隐藏状态的递归传递，使当前时刻的输出不仅依赖当前输入，还依赖以往历史信息，实现了对时间上下文的自然建模，并凭借参数共享机制在理论上可以处理任意长度的序列，因此在早期的时间序列预测、语音建模和自然语言处理等任务中得到广泛应用。 然而，基础 RNN 在实践中暴露出明显短板：一是梯度在时间反向传播过程中易出现消失或爆炸，导致模型难以有效学习远距离时间步之间的长程依赖；二是串行计算特性使其难以充分利用现代并行硬件，训练效率受限。为缓解这些问题，长短期记忆网络（Long Short-Term Memory, LSTM）和门控循环单元（Gated Recurrent Unit, GRU）相继被提出，通过在循环结构中引入记忆单元与门控机制，有选择地保留或遗忘历史信息，显著缓解了梯度消失问题，使得基于 RNN/LSTM/GRU 的模型在较长时间跨度的预测任务中表现出更好的稳定性和精度\cite{cho2014learning}。在此基础上，大量变体陆续提出，例如通过膨胀连接扩大感受野的 Dilated RNN\cite{chang2017dilated}、通过双阶段注意力机制自适应选择关键时间步与特征的 DA-RNN，以及结合序列到序列框架与分位数回归以给出区间预测的 MQ-RNN 等\cite{qin2017dual,wen2017multi}，这些方法进一步提升了循环结构在长序列、带不确定性的预测任务中的适用性。 与 RNN 系列侧重沿时间轴递归建模不同，卷积神经网络（Convolutional Neural Network, CNN）最初主要服务于图像等空间数据，通过局部感受野与权值共享机制高效提取局部模式。随后研究者将一维卷积推广到时间轴上，发展出适用于时间序列的 1D CNN 结构，通过滑动卷积核对局部时间窗口进行特征提取，从而挖掘短期模式与局部依赖关系。代表性工作如 WaveNet 利用一维膨胀卷积在语音信号建模中显著提升了对长程依赖的捕捉能力，展示了基于卷积的时序建模在长序列任务中的潜力\cite{oord2016wavenet}。在此基础上，Temporal Convolutional Network（TCN）通过多层一维卷积堆叠与膨胀因子的逐层扩展，在保持严格因果性的前提下获得了大感受野，兼具并行计算优势与长程建模能力，在时间序列预测、信号处理和序列建模等任务中均取得了优异表现\cite{bai2018empirical}。 在循环结构与卷积结构各自发展的同时，学界也开始探索 RNN 与 CNN 的融合架构，以兼顾局部模式提取与长程依赖建模。典型做法是首先利用 CNN 在时间维或时空图上抽取局部/空间特征，再将提取到的高层次特征序列输入 LSTM/GRU 等循环单元进行时间建模。例如，面向交通流量预测的 DCRNN 将交通网络建模为有向图，利用扩散卷积刻画空间依赖关系，再结合门控循环结构捕捉时间动态，从而在时空依赖建模方面取得显著优势；TPA-LSTM 则通过卷积模块提取局部时间模式，并结合 LSTM 与注意力机制学习关键时间片段，提高了在复杂周期性和多尺度波动场景下的预测精度\cite{li2018diffusion,shih2019temporal}。综合来看，基于 RNN 与 CNN 的方法在较长一段时间内构成了深度学习时间序列预测的主流方案，相较于传统统计与浅层机器学习方法，显著提升了对非线性关系、长程依赖和多尺度结构的建模能力，为后续基于注意力机制和 Transformer 的模型奠定了重要基础。 \subsection{基于 Transformer 与注意力机制的方法} 在循环与卷积结构推动时间序列预测取得显著进展的基础上，Transformer 的提出进一步改变了序列建模的范式。作为一种以自注意力机制为核心的编码--解码架构，Transformer 最初在机器翻译等自然语言处理任务中取得突破性进展。其通过多头自注意力在全局范围内建模序列各位置间的依赖关系，并依托完全并行的计算模式显著提升了长序列建模的效率与表达能力。与递归结构依赖时间步串行计算不同，自注意力可以在单层中同时访问完整的上下文信息，根据相关性为不同时间步分配不同权重，从而在理论上更适合捕捉长程依赖与复杂模式，这使得 Transformer 很快被引入时间序列预测领域\cite{chen2025transformercn}。 在时间序列预测，尤其是长程时间序列预测（Long-Term Series Forecasting, LTSF）任务中，基于 Transformer 的模型近年来成为研究热点。一方面，自注意力机制在时间域上可同时考虑所有历史时间步，并通过学习的注意力权重自动突出关键时间片，有助于挖掘多尺度趋势与周期模式；另一方面，多头结构与残差连接等设计增强了模型对高维、多变量时序数据的表示能力。然而，原始 Transformer 直接应用于长序列时会面临自注意力计算与存储复杂度随序列长度呈二次增长的问题，且自回归式解码容易导致误差在多步预测中累积。为缓解这些问题，大量改进型 Transformer 被提出：一类工作通过稀疏注意力、局部窗口或金字塔结构减少注意力连接数，从而将复杂度从平方级降低到近似 L log L 量级。其中，LogTrans 通过日志稀疏模式压缩远距离依赖连接\cite{li2019enhancing}；Reformer 引入局部敏感哈希与可逆残差结构以降低计算和存储开销\cite{kitaev2020reformer}；Informer 则借助 ProbSparse 注意力与生成式解码机制提升长序列预测效率。另一类工作则从时间序列特性出发，将趋势--季节分解、频域表示等先验融入 Transformer 结构，如 Autoformer 将序列分解与自相关机制内嵌于网络模块，Pyraformer 采用金字塔式注意力以分层方式建模长依赖\cite{liu2022pyraformer}，FEDformer 则利用 Fourier/Wavelet 稀疏基在频域中压缩表示长程周期信息。这些模型在多个公共数据集上相较基础 Transformer 与传统深度模型取得了更优的长程预测性能。 尽管上述基于注意力机制的模型在长程依赖建模和性能上取得了重要进展，Transformer 在时间序列预测中的应用仍存在若干共性瓶颈。一方面，即便采用稀疏或分块注意力，其核心计算模式仍然显著重于线性模型或卷积模型，难以彻底消除随序列长度增长带来的时间与空间开销，在资源受限或实时性要求较高的场景中部署成本较高。另一方面，近期研究表明，简单增加时间窗口长度并不总能提升 Transformer 在长程预测任务中的表现，反而可能导致模型在过长输入上过拟合噪声，难以真正提取有效的远程时序信息。此外，归一化与残差等设计在处理非平稳时间序列时可能削弱关键幅值信息，促使后续工作提出去平稳化注意力\cite{liu2022nonstationary}和频域增强等改进方向，以提升 Transformer 在复杂时序分布下的稳健性。总体来看，基于 Transformer 与自注意力机制的方法在时间序列预测领域已经形成丰富的模型家族，为长程依赖建模提供了强大工具，但在效率、鲁棒性以及与其他结构（如线性模型、频域网络等）的协同方面仍存在改进空间，也为本文后续工作提供了重要的研究基础与对比对象。

\subsection{基于大语言模型(LLM)的方法} 近年来，大语言模型（Large Language Model, LLM）在自然语言处理和多模态理解领域取得了突破性进展，研究者开始尝试将其强大的表征与推理能力迁移到时间序列分析中，用于预测、分类、异常检测、填补缺失等多种任务。在电力等典型垂直场景中，国内综述研究也表明 LLM 在少样本迁移、多模态信息融合和复杂预测情境适应方面具有较大潜力\cite{zhang2025llmpowercn}。国际上的综述与观点研究则进一步指出，基于 LLM 的时间序列方法正在从零样本推理逐步拓展到提示学习、离散化建模、跨模态对齐与统一建模等更丰富的技术路线\cite{zhang2024llmtsurvey,jin2024position}。与前述仅针对数值序列设计的深度模型不同，LLM 方法的核心挑战在于如何弥合离散文本预训练与连续数值时间序列之间的模态鸿沟，并在此基础上有效地复用预训练模型中蕴含的大规模先验知识。围绕这一目标，现有工作大致可以分为五类路径：直接提示、量化建模、模态对齐、视觉桥接以及工具集成。

第一类是基于提示（prompting）的直接调用方法，将时间序列以文本形式序列化后直接输入现成的 LLM，通过设计合适的提示模板实现零样本或小样本的预测与分类。例如，将温度、电价或传感器读数按时间顺序写入自然语言描述，再让模型输出未来趋势或类别；还有工作引入语义空间感知的提示学习，以增强时序表示与语言先验之间的联系\cite{gruver2024zeroshot,pan2024s2ipllm}。这类方法实现成本低、无需额外训练，适合缺乏标注数据的场景，但将精细数值信息转成长文本既拉长上下文，又容易丢失数值间的精确关系，在高维、多变量或高精度预测任务中效率与效果均受限。

第二类是基于量化（quantization）的离散化方法，通过向量量化（如 VQ-VAE）、K-Means 聚类或离散分箱等手段，将连续时间序列映射到有限的代码本索引或符号类别，再作为 token 输入 LLM。代表性工作如 TOTEM 从跨域通用离散表示的角度探索时间序列 token 化，使模型能够以更接近语言建模的方式处理不同任务和不同来源的时序数据\cite{talukder2024totem}。这种离散表示既让时间序列更贴近 LLM 原生的离散输入形式，便于复用现有架构，也可以通过在时间或频率域上降采样与编码显著压缩长序列长度，缓解上下文窗口和计算复杂度压力。然而，这类方法通常需要额外训练编码器/解码器，存在两阶段训练带来的子最优问题，同时量化误差和代码本设计也会影响下游任务性能。

第三类是基于模态对齐（aligning）的方法，通过显式学习时间序列与语言空间之间的一致性来迁移 LLM 能力。一部分工作采用对比学习、最优传输等目标，将时间序列编码器输出与文本描述或报告的嵌入对齐，学习共享语义空间，从而支持跨模态检索、报告生成等任务；另一部分工作则以预训练 LLM（如 GPT 系列、LLaMA 系列）为主干，在其输入前叠加专门的时间序列编码层，将数值序列投影到适配 LLM 的嵌入空间，并通过微调或适配器实现统一处理多种时序任务（预测、分类、异常检测等）。这一路线已经出现了较多代表性工作，如 LLM4TS、TimeCMA 与 UniTime 等\cite{chang2025llm4ts,liu2025timecma,liu2024unitime}。其优势在于能够端到端地保留 LLM 的语言知识和推理能力，但模型设计与训练过程相对复杂，对计算和数据资源要求较高。

第四类是以视觉模态为桥梁的多模态方法，将时间序列转化为图像或与视觉信号配对，再利用预训练视觉--语言模型或多模态 LLM 完成对时间序列的理解与生成。例如，把时间序列绘制为折线图、频谱图等输入视觉--语言模型，让其生成趋势描述、风险提示等文本，或者通过多模态预训练将 IMU、EEG 等时序信号对齐到图像--文本共享空间。代表性工作如 Time-VLM 将时间序列图像化并结合视觉--语言模型进行增强预测，展示了视觉桥接路线在少样本与零样本场景中的潜力\cite{zhong2025timevlm}。这类方法能够间接利用成熟的视觉多模态模型，但额外引入了图像渲染和视觉编码过程，难以精确保留原始数值信息，也并非适用于所有类型的时间序列数据。

第五类是以 LLM 作为外部工具调度器（tool integration）的方法，并不直接让 LLM 对时间序列做前向预测，而是让其根据文本指令或先验知识生成代码、配置或 API 调用，再由专用时序模型执行具体分析。例如，利用 LLM 自动生成损失函数、特征工程代码或调用天气、金融等外部时序 API 进行推理。此类方法充分发挥了 LLM 在自然语言理解与程序生成方面的优势，为时间序列任务提供元建模与自动化建模的能力，但整体流程往往难以端到端联合优化。

基于大语言模型的时间序列方法已经展示了在多任务、少样本甚至零样本场景下的潜力，能够将丰富的语言与领域知识引入数值时序建模，但也面临数值表示不稳定、上下文长度与计算开销受限、模态对齐仍不充分等挑战。如何在保持 LLM 语义与推理优势的同时，设计更高效的数值表示与跨模态对齐机制，并结合频域分析、时序分解等领域知识构建统一的时序--语言建模框架，已成为当前研究的重点方向，也为本文后续基于大语言模型的多元时间序列预测框架提供了重要的理论与方法基础。
\section{本文研究目标与主要内容}

针对多元长程时间序列预测中精度、长程建模能力与结构复杂度之间难以平衡的问题，本文围绕统一预测骨架构建，局部多尺度频域增强和全局频域增强三条主线开展研究，目标是在统一骨架下提升模型对长期趋势、多尺度波动和跨变量关联的刻画能力。全文的主要工作可概括为以下三个方面：

（1）构建统一预测骨架。以时域编码为主体，引入外部冻结大语言模型生成的语义嵌入和基础频域表示，通过轻量适配器、步长感知融合和跨模态对齐实现多源信息协同建模，为后续不同频域增强模块的替换提供统一接口。

（2）提出融合小波包分解的局部多尺度频域增强方法。通过多尺度子带分解和子带级特征融合，增强模型对局部突变和复杂周期结构的刻画能力，并验证局部多尺度频域增强对复杂局部波动与高误差样本建模的有效性。

（3）提出融合频域前馈网络的全局频域增强方法。通过复数频谱建模、稀疏化约束和门控融合提升全局周期结构建模能力，并考察其在综合预测性能、长期趋势刻画和平均误差控制方面的表现。

本文尝试构建一套兼顾预测精度与结构可解释性的多元时间序列预测方案，为大模型在时序建模中的应用提供系统化研究路径，并揭示局部多尺度频域建模与全局频域建模在作用侧重点上的差异。

\section{论文组织结构}
全文共分为六章。第一章为绪论，介绍研究背景、国内外研究现状以及本文的研究目标。第二章阐述时间序列预测、大语言模型、频域分析和门控注意力等相关理论基础。第三章提出基于 Qwen3 语义先验的统一预测骨架，并验证其有效性。第四章研究融合小波包分解的局部多尺度频域增强方法。第五章研究融合频域前馈网络的全局频域增强方法。第六章对全文工作进行总结，同时讨论当前研究局限与未来发展方向。
